\documentclass[11pt,a4paper]{article}
\usepackage[margin=24mm]{geometry}
\usepackage{amsmath,amssymb}
\usepackage[hidelinks]{hyperref}
\newcommand{\appendices}{\appendix}
\title{Paper 1: Technical Revision Fragments}
\author{Review copy for Abdullah Haatim}
\date{7 October 2026}
\begin{document}
\maketitle
\noindent This note previews insertable revisions for \textit{ThesisProj-7.pdf}.
It is not the full manuscript and introduces no new experimental results.
The exact insertion map and source qualification are in the accompanying
\texttt{README.txt}. The section numbering below belongs to this preview.
\section{Related-work insertion}

% Insert in the Introduction immediately before "We propose a
% Hankel-structured extension of EM-GS...". Keep the existing introduction.

Low-rank Hankel modeling and Cadzow processing are established tools
\cite{cadzow1988,markovsky2008}. Related uses of structured low-rank
matrices include spectral compressed sensing and MRI reconstruction
\cite{chen2014emac,jin2016aloha,haldar2014loraks}. Our contribution is not
a new general low-rank recovery principle. We study how a simple spatial
Hankel correction behaves inside a magnitude-only Rydberg channel
estimator, including conditions where it gives little benefit or causes
harm. The Hankel matrix is formed along each user's array dimension.
This construction differs from a rank constraint on a two-dimensional
channel slice, as considered in~\cite{xu2025}.

Hard truncation gives a best rank-limited approximation to its input
matrix in the Frobenius norm. It need not give the best estimate of the
underlying clean channel. Singular-value shrinkage methods such as
OptShrink address this distinction~\cite{nadakuditi2014optshrink}. Their
matrix-noise assumptions require care here: Hankel lifting repeats
entries, and the error entering the structural step comes from an
iterative channel estimate. We therefore do not assume that a shrinkage
rule is optimal for this problem simply because it is optimal under a
different denoising model.


\section{Replacement estimator text}

% Replace all of Section III-B, "The Estimator".
\subsection{The Estimator}
\label{sec:revised-estimator}

HS-GS uses the same measurement update and spectral initialization as
EM-GS. It then applies the Hankel step separately to each user column
$\mathbf g_k$. We do not stack different users or pilot measurements into
one Hankel matrix, and we do not assume that users share arrival angles.

At iteration $t$, the measurement update is
\begin{align}
\mathbf M^{(t)}&=\mathbf G^{(t)}\mathbf S+\mathbf B,\\
\boldsymbol\kappa^{(t)}&=
  2\mathbf Z\odot|\mathbf M^{(t)}|/\sigma^2,\\
\mathbf Y^{(t)}&=\mathbf Z\odot R(\boldsymbol\kappa^{(t)})
  \odot e^{\jmath\arg\mathbf M^{(t)}},\\
\widetilde{\mathbf G}^{(t+1)}(\mathbf S\mathbf S^H)
  &=(\mathbf Y^{(t)}-\mathbf B)\mathbf S^H,
\label{eq:revision-em-update}
\end{align}
where $R(x)=I_1(x)/I_0(x)$ and all nonlinear operations are elementwise.
The last equation is solved as a linear system. The pilot matrix for each
fit is assumed to have full row rank. Scaled Bessel functions are used to
evaluate the ratio stably; the contribution to $\mathbf Y^{(t)}$ is zero
when the corresponding entry of $\mathbf M^{(t)}$ is zero.

For a channel vector of length $N$, we use
\begin{equation}
p=\left\lfloor\frac{N-1}{2}\right\rfloor,
\qquad
[\mathcal H_p(\mathbf g)]_{a+1,b+1}=g[a+b+1],
\end{equation}
with $0\leq a<N-p$ and $0\leq b\leq p$. This gives an
$(N-p)\times(p+1)$ matrix and the exact capacity
$r_{\max}=\lceil N/2\rceil$.

One Cadzow sweep is
\begin{equation}
\mathcal P_r(\mathbf g)=
\mathcal H_p^\dagger\!\left(\Pi_r(\mathcal H_p(\mathbf g))\right).
\label{eq:revision-cadzow}
\end{equation}
Here $\Pi_r$ keeps the $r$ largest singular values and sets the others to
zero. It does not change the retained singular values. The operation
$\mathcal H_p^\dagger$ averages the entries on each antidiagonal to recover
the channel vector. We apply four sweeps to every column of
$\widetilde{\mathbf G}^{(t+1)}$ before the next measurement update.

This is a finite sequence of structural corrections. After antidiagonal
averaging, the Hankel matrix formed again from the output can have rank
greater than $r$. Low Hankel rank also does not impose every physical
constraint in the geometric channel model. A small rank can remove useful
channel components, while a large rank may leave much of the estimation
error unchanged.

% Replace all of Section III-C, "Choosing the Hankel Rank".
\subsection{Choosing the Hankel Rank}
\label{sec:revised-rank-selection}

We choose the rank from the available pilots without using the true path
count. The last $P_B=\max\{1,\operatorname{round}(0.3P)\}$ pilot columns
form the validation set $B$. The remaining columns form the fitting set
$A$. For example, with $P=30$, we use 21 columns to fit each candidate and
nine columns to compare the candidates.

For every $r\in\{1,\ldots,r_{\max}\}$, we run the estimator on set $A$
and obtain $\widehat{\mathbf G}^{(r)}_A$. Its initialization also uses only
set $A$. We then select
\begin{equation}
\widehat r=\operatorname*{arg\,min}_{1\leq r\leq r_{\max}}
\left\|\mathbf Z_B-
 \left|\widehat{\mathbf G}^{(r)}_A\mathbf S_B+\mathbf B_B\right|
\right\|_F^2.
\label{eq:revision-rank-selection}
\end{equation}
The score sums the squared magnitude residuals over all receive elements
and validation columns. If scores are exactly tied, we select the smallest
rank. One scalar rank is shared by all user columns; the Hankel operations
on those columns remain separate.

After choosing the rank, we restart the estimator from an initialization
computed using all $P$ pilots and fit the final estimate using all of them.
Thus, the validation split does not permanently discard 30\% of the pilots
or require additional pilot transmissions. Its costs are the repeated
candidate fits and the smaller fitting set used during rank selection.

If $\widehat r=r_{\max}$, the structural step is bypassed. With the same
initialization and update count, the final estimate then equals the paired
EM-GS estimate exactly. This fallback does not guarantee that a smaller
selected rank improves the unknown channel error in every trial.

\begin{center}
\fbox{\begin{minipage}{0.91\columnwidth}
\small
\textbf{HS-GS procedure}\par
\textit{Inputs:} $\mathbf Z,\mathbf S,\mathbf B,\sigma^2$; candidate and
final update counts $T_{\rm sel},T$; four Cadzow sweeps per application.
\begin{enumerate}
\setlength{\itemsep}{1pt}
\setlength{\parskip}{0pt}
\item Form fitting set $A$ and validation set $B$ from the pilot columns.
\item For each candidate rank, initialize from $A$ and perform
$T_{\rm sel}$ measurement updates, each followed by four per-user Cadzow
sweeps. Bypass these sweeps at $r_{\max}$.
\item Select $\widehat r$ using \eqref{eq:revision-rank-selection}.
\item Reinitialize using all pilots. Perform $T$ measurement updates with
the selected structural processing and return the final channel estimate.
\end{enumerate}
\end{minipage}}
\end{center}


\section{Replacement bound text}

% Replace all of Section IV-B, "Constrained Cramer--Rao Bound".
\subsection{Constrained Cram\'er--Rao Bound}
\label{sec:revised-ccrb}

The ordinary CRLB treats the real and imaginary parts of each channel
coefficient as separate unknowns. The geometric CCRB also uses the fact
that these coefficients are generated by a common set of paths. It is
\begin{equation}
\mathbf C_{\rm geom}
=\mathbf U(\mathbf U^T\mathbf J\mathbf U)^{-1}\mathbf U^T,
\label{eq:revision-ccrb}
\end{equation}
where $\mathbf J$ is the Fisher information matrix in real channel
coordinates. The columns of $\mathbf U$ form an orthonormal basis for the
local changes in the channel allowed by the geometric model
\cite{gorman1990,stoica1998}. Appendix~\ref{app:revision-ccrb} gives the
parameter ordering, derivatives and numerical calculation.

The bound is evaluated at the true channel and its path configuration.
The path gains and spatial frequencies are still unknown coordinates in
the calculation; they are not supplied to the estimator. The known
quantities are the pilots, reference field and noise variance. The true
model order is used only to define this ideal geometric reference.

We plot the trace of each bound with the same channel-energy normalization
as the NMSE. The gap between the two references shows the potential value
of the geometric information under their local unbiasedness assumptions.
It does not prove that the gap is achievable by HS-GS or another practical
estimator. EM-GS and HS-GS are not assumed to be unbiased, so neither curve
is treated as an absolute lower limit on their measured NMSE.


\section{Short scope clarifications}
\subsection{Effective-rank aggregation}

% Insert in Section V-E, "Effective Rank", immediately after the definition
% of rho and before "The value rho can be understood as the fraction...".

We calculate this quantity separately for each user column. At an
operating point, the reported $\rho$ is the median over the user columns
and channel realizations. When resampling trials, we keep all users from
the same realization together. The values shown here are calculated from
the true noiseless channels and are used to explain the simulation results;
they are not inputs to rank selection.


\subsection{Interpretation of the existing bound gaps}

% In Section V-B, replace the two paragraphs starting "The geometric CCRB
% is approximately..." and "HS-GS uses part of this structure...".
% Keep the numerical values and the following unbiasedness caveat.

The geometric CCRB is approximately $6.95$--$7.06$~dB below the ordinary
CRLB. It includes the local geometric multipath structure, while the
ordinary CRLB treats the channel coefficients as separate unknowns.
HS-GS improves on EM-GS but remains $4.37$--$4.85$~dB above the geometric
reference. These gaps describe the two reference curves and the measured
estimator performance under this model. They do not establish how much
further improvement a practical estimator can achieve.


\subsection{Existing projection-placement comparison}

% Insert at the end of the first paragraph of Section V-G,
% "Projection Placement and Computation", after the sentence mentioning
% 300 paired trials and four Cadzow sweeps.

The common rank is selected using the interleaved candidates. Thus, this
comparison measures placement at a shared selected rank; it does not tune
each variant separately. The four sweeps are used at each structural
application, so interleaving and final-only processing have different
total computational costs.

% Insert after the conclusion and before the bibliography.
% If the manuscript already starts its appendices elsewhere, omit the next line.
\appendices
\section{Real-Coordinate Fisher Information and Geometric CCRB}
\label{app:revision-ccrb}

\subsection{Channel coordinates and measurement information}

For receive element $n$, define
\begin{equation}
\mathbf x_n=
 [\Re G_{n1},\ldots,\Re G_{nK},
  \Im G_{n1},\ldots,\Im G_{nK}]^T.
\end{equation}
The full real vector is
$\mathbf x=[\mathbf x_1^T,\ldots,\mathbf x_N^T]^T\in\mathbb R^{2NK}$.
This ordering gives
$\|\widehat{\mathbf x}-\mathbf x\|_2^2
=\|\widehat{\mathbf G}-\mathbf G\|_F^2$.
We condition on the known $\mathbf S$, $\mathbf B$ and $\sigma^2$.
The noise convention is $\mathbb E|W_{np}|^2=\sigma^2$, so each real
component has variance $\sigma^2/2$.

Write $\mu_{np}=(\mathbf G\mathbf S+\mathbf B)_{np}$ and
$a_{np}=|\mu_{np}|$. For $z\geq0$, the amplitude density is
\begin{equation}
p(z\mid a)=\frac{2z}{\sigma^2}
 \exp\!\left(-\frac{z^2+a^2}{\sigma^2}\right)
 I_0\!\left(\frac{2za}{\sigma^2}\right).
\end{equation}
With $R(u)=I_1(u)/I_0(u)$, its amplitude score is
\begin{equation}
\frac{\partial\log p(z\mid a)}{\partial a}
=\frac{2}{\sigma^2}
 \left[zR\!\left(\frac{2za}{\sigma^2}\right)-a\right].
\end{equation}
Using $\mathbb E[zR(2za/\sigma^2)]=a$, the scalar information is
\begin{align}
\mathcal I_a(a)
&=\frac{4}{\sigma^4}
 \left[\mathbb E\!\left\{z^2R^2\!\left(\frac{2za}{\sigma^2}\right)\right\}
 -a^2\right],
\label{eq:revision-scalar-information}\\
\beta(a)&=\mathcal I_a(a)/4.
\end{align}

For $\mu_{np}\ne0$, put
\begin{equation}
\mathbf c_{np}=e^{-\jmath\arg\mu_{np}}\mathbf S_{:,p},
\qquad
\mathbf v_{np}=
\begin{bmatrix}\Re\mathbf c_{np}\\-\Im\mathbf c_{np}\end{bmatrix}.
\end{equation}
Here $\mathbf v_{np}=\partial a_{np}/\partial\mathbf x_n$.
Independence of the measurement noises gives
\begin{align}
\mathbf J_n&=\sum_{p=1}^{P}
  \mathcal I_a(a_{np})\mathbf v_{np}\mathbf v_{np}^T,\\
\mathbf J&=\operatorname{blkdiag}
 (\mathbf J_1,\ldots,\mathbf J_N).
\label{eq:revision-real-fisher}
\end{align}
Each nonzero measurement term has rank one in real coordinates. At
$\mu_{np}=0$, the contribution is defined by its continuous zero limit;
an arbitrary choice of phase supplies no information.

\subsection{Geometric tangent space}

For each user $k$ and path $\ell$, take the real parameters
$(\psi_{\ell,k},\Re\alpha_{\ell,k},\Im\alpha_{\ell,k})$.
Collect them in $\boldsymbol\eta$ and set
$\mathbf D=\partial\mathbf x/\partial\boldsymbol\eta$.
Writing $v_n(\psi)=e^{-\jmath(n-1)\psi}$, the complex derivatives for
one path are
\begin{align}
\frac{\partial g_k[n]}{\partial\psi_{\ell,k}}
 &=-\jmath(n-1)\alpha_{\ell,k}v_n(\psi_{\ell,k}),\\
\frac{\partial g_k[n]}{\partial\Re\alpha_{\ell,k}}
 &=v_n(\psi_{\ell,k}),\\
\frac{\partial g_k[n]}{\partial\Im\alpha_{\ell,k}}
 &=\jmath v_n(\psi_{\ell,k}).
\end{align}
Their real and imaginary parts fill the corresponding entries of
$\mathbf x_n$; entries belonging to other users are zero. Using spatial
frequency $\psi$ follows the channel generator directly and avoids a
coordinate singularity from $\psi=\pi\sin\theta$ at endfire.

Take a thin SVD $\mathbf D=\mathbf U_D\boldsymbol\Sigma_D\mathbf V_D^T$.
The implementation retains the columns of $\mathbf U_D$ whose singular
values exceed $10^{-9}$ times the largest singular value. These columns
form $\mathbf U$, an orthonormal basis for the numerically resolved range
of $\mathbf D$. This construction avoids inverting a redundant physical
parameterization when $\mathbf D$ is not full column rank.

For locally unbiased estimation along this tangent space,
\begin{equation}
\mathbf C_{\rm geom}
=\mathbf U(\mathbf U^T\mathbf J\mathbf U)^{-1}\mathbf U^T.
\end{equation}
The gains and spatial frequencies are varied jointly in $\mathbf D$;
neither is held fixed as known side information. This calculation includes
no unknown reference-field, calibration or noise-variance parameters.
Adding such nuisance parameters would require the corresponding enlarged
information matrix. The true path counts and parameter point specify the
local reference, rather than inputs available to HS-GS.

The displayed inverses require positive definite information on the
relevant space. A singular restricted matrix indicates an unresolved
local direction and is not replaced by a regularized finite bound.
The known reference can remove phase ambiguity, but its presence alone
does not guarantee identifiability for every pilot design. Coincident or
zero-gain paths can also invalidate a regular local parameterization.
The stated numerical rank tolerance must therefore accompany the results.

\subsection{Numerical evaluation and normalization}

For trial $i$, the two scalar reference errors are
\begin{align}
b_{{\rm unc},i}&=\sum_{n=1}^{N}\operatorname{tr}(\mathbf J_{n,i}^{-1}),\\
b_{{\rm geom},i}&=
 \operatorname{tr}\!\left[(\mathbf U_i^T\mathbf J_i\mathbf U_i)^{-1}\right].
\end{align}
The second identity uses $\mathbf U_i^T\mathbf U_i=\mathbf I$.
For either reference, the plotted value is
\begin{equation}
10\log_{10}
 \frac{\sum_i b_i}{\sum_i\|\mathbf G_i\|_F^2}.
\label{eq:revision-bound-normalization}
\end{equation}
The same channel realizations and energy denominators are used for the
estimator curves. This is a ratio of sums, not an average of trial-wise
dB values.

The scalar expectation is evaluated with scaled Bessel functions and
quadrature over $[\max(0,a-8\sigma),a+8\sigma]$, where
$\sigma=\sqrt{\sigma^2}$. The quadrature uses absolute tolerance
$10^{-14}$ and relative tolerance $10^{-12}$, splitting at $a$ when it
lies inside the integration window. At each main-figure SNR, 800 equally
spaced amplitude nodes cover $0.99a_{\min}$ to $1.01a_{\max}$ over that
SNR's worlds. Piecewise cubic Hermite interpolation supplies the scalar
weights at the measurement amplitudes. The implementation checks 200
sampled amplitudes against direct quadrature and requires maximum relative
error below $10^{-4}$. The tangent derivatives are checked by finite
differences. It also checks that the constrained trace is positive and
no larger than the unconstrained trace, and records the condition number
of the restricted information matrix. No diagonal ridge is added.


\begin{thebibliography}{9}

\bibitem{xu2025}B.~Xu, J.~Zhang, Z.~Chen, B.~Cheng, Z.~Liu, Y.-C.~Wu, and B.~Ai,
``Channel estimation for Rydberg atomic receivers,''
\emph{IEEE Wireless Commun. Lett.},
vol.~14, no.~9, pp.~2957--2961, 2025.

\bibitem{chen2014emac}Y.~Chen and Y.~Chi,
``Robust spectral compressed sensing via structured matrix completion,''
\emph{IEEE Trans. Inf. Theory},
vol.~60, no.~10, pp.~6576--6601, 2014.

\bibitem{jin2016aloha}K.~H.~Jin, D.~Lee, and J.~C.~Ye,
``A general framework for compressed sensing and parallel MRI using annihilating filter based low-rank Hankel matrix,''
\emph{IEEE Trans. Comput. Imag.},
vol.~2, no.~4, pp.~480--495, 2016.

\bibitem{haldar2014loraks}J.~P.~Haldar,
``Low-rank modeling of local $k$-space neighborhoods (LORAKS) for constrained MRI,''
\emph{IEEE Trans. Med. Imag.},
vol.~33, no.~3, pp.~668--681, 2014.

\bibitem{cadzow1988}J.~A.~Cadzow,
``Signal enhancement---a composite property mapping algorithm,''
\emph{IEEE Trans. Acoust., Speech, Signal Process.},
vol.~36, no.~1, pp.~49--62, 1988.

\bibitem{markovsky2008}I.~Markovsky,
``Structured low-rank approximation and its applications,''
\emph{Automatica},
vol.~44, no.~4, pp.~891--909, 2008.

\bibitem{gorman1990}J.~D.~Gorman and A.~O.~Hero,
``Lower bounds for parametric estimation with constraints,''
\emph{IEEE Trans. Inf. Theory},
vol.~36, no.~6, pp.~1285--1301, 1990.

\bibitem{stoica1998}P.~Stoica and B.~C.~Ng,
``On the Cram\'er--Rao bound under parametric constraints,''
\emph{IEEE Signal Process. Lett.},
vol.~5, no.~7, pp.~177--179, 1998.

% Add this entry inside the existing thebibliography environment.
% The DOI identifies the IEEE publication. This uses no unverified page range.
\bibitem{nadakuditi2014optshrink}
R.~R.~Nadakuditi,
``OptShrink: An algorithm for improved low-rank signal matrix denoising by
optimal, data-driven singular value shrinkage,''
\emph{IEEE Trans. Inf. Theory}, 2014,
doi: 10.1109/TIT.2014.2311661.


\end{thebibliography}
\end{document}
